{\large\textbf{E: Physics of illumination}}

An incandescent light bulb produces light by heating a tungsten filament to a
high enough temperature to emit black body radiation in the visible part of the
spectrum, but a lot of energy is still wasted in the infrared part.

To quantify light as perceived by human vision, we use photometric units, which
take into account sensitivity of the human eye to light of different
wavelengths. The total amount of visible light \textit{emitted} by a source into
all directions is called \textbf{luminous flux}, measured in lumens
$[\mathrm{lm}]$. The amount of visible light \textit{received} by a surface per
unit of its area is called \textbf{illuminance}, in lux
$[\mathrm{lx} = \mathrm{lm/m^2}]$, and can be measured by a light meter.

When measuring quantities of light based solely on the energy it carries, we
use radiometric units, which are expressed with conventional units of power.
Radiometric counterpart of the luminous flux is \textbf{radiant flux}, measured
in watts $[\mathrm{W}]$, and \textbf{irradiance} $[\mathrm{W/m^2}]$ is the
counterpart of the illuminance.

Today you will study thermal and illumination properties of light sources. The
three tasks are mostly independent from each other. Sketch your setup for each
task. No error analysis is required for Tasks 1 \& 2.

\textbf{Equipment (see also Fig. 1)}

A\quad Black and white plastic plate $3\,\mathrm{mm}$ thick with a stand. Both
plates are good absorbers of infrared light.

B\quad Light meter with a stand. The light meter turns off automatically after
6 minutes -- turn it back on with a long press on the on/off button. Pay
attention to the units (lx, not fc). You may use the \texttt{HOLD} button to
freeze the displayed value.

C\quad A light mounting stand with a round base, a weight for stability, and two
interchangeable light modules: incandescent light bulb (\textbf{maximum
voltage} $12\,\mathrm{V}$) and LED (\textbf{maximum voltage}
$3.0\,\mathrm{V}$, \textbf{do not exceed} $400\,\mathrm{mA}$
\textbf{current}). You may use toothpicks to wedge the modules in place. Black
paper is provided to shield your eyes while reading the instruments.

D\quad Infrared thermometer. The measurement reading is held after a short
delay when the trigger is \textit{released}. The measurements may have a
substantial but constant systematic error.

E\quad Paper work mat with angular and distance grid.

F\quad Protractor.

G\quad Red, green and blue light filters in an envelope. If you have trouble
telling the colours apart, raise the help card for assistance.

\noindent\fbox{\parbox{\dimexpr\linewidth-2\fboxsep-2\fboxrule\relax}{\textbf{The
filters are sensitive to heat. Keep them away from the light source.}}}

H\quad Power supply. Press the voltage/current knob \textit{multiple times} to
select which digit to adjust (indicated by the blinking light below the digit),
and turn the knob to change the digit. After a few seconds, the light stops
blinking and the display starts showing the actual voltage/current. Vary the
current to control the light source. If the requested current cannot be reached
without exceeding the maximum voltage, the power source will switch into the
constant voltage mode and limit the current. Plug the wires into the matching
negative (black) and positive (red) sockets of the power supply. Do not use the
green socket.

\noindent\fbox{\parbox{\dimexpr\linewidth-2\fboxsep-2\fboxrule\relax}{\textbf{To
avoid damaging the light sources, set the voltage to allowed maximum and set the
current to zero before you plug in the wires!}\\
\textbf{If your light source burns out, you can ask for a replacement. Note that
only a limited number of spare light sources are available.}}}

\textbf{Task 1 - Colour and temperature (4 pts)}

The colour of the black body radiation depends on its temperature. In
astronomy, the temperature of stars is determined from their colour index, the
ratio of illuminances measured through two different colour filters.

\textbf{(a)} Table 1 contains the illuminances measured through the red, green
and blue filter for a standard incandescent light source at known temperatures.
Choose suitable light filters and construct a calibration curve that relates
the chosen colour index to the temperature.

\textbf{(b)} Measure the relationship between the electrical input power and
the tungsten filament temperature. Plot the result over a relevant range.

\textbf{Task 2 - Luminous efficacy (8 pts)}

The performance of light sources is quantified by their \textbf{luminous
efficacy}, measured in lumens per watt, as the ratio between the luminous flux
and the consumed power. As a point of reference, the sun has luminous efficacy
of $93\,\mathrm{lm/W}$.

Measure the dependence of luminous efficacy on the electrical input power for
both light sources across the range with detectable light output. Plot the
results, one plot per light source. State all steps of the calculation
procedure and present all the measured data.

\textbf{Task 3 - Radiative heating (8 pts)}

\noindent\fbox{\parbox{\dimexpr\linewidth-2\fboxsep-2\fboxrule\relax}{\textbf{The
following task may be time consuming, plan your work accordingly.}}}

When light hits an object, some of it is absorbed. At moderate temperature
differences between the object and the environment, we can model heat
dissipation into the surroundings with the \textbf{heat transfer coefficient}
$h$, in the form $P/A = h(T - T_0)$, where $T$ is the temperature of the
surface, $T_0$ the temperature of the surroundings, and $P/A$ denotes the power
lost to the environment due to dissipation, per unit area.

\textbf{(a)} Determine the heat transfer coefficient $h$ and the thermal
conductivity $\lambda$ for the black plastic, and perform error analysis.
Assume the material absorbs all received light and the incandescent light bulb
emits all power in the form of electromagnetic radiation.

\textbf{(b)} Estimate the albedo (the fraction of the irradiance that is
reflected instead of absorbed) of the white plastic and perform error analysis.

\textbf{Useful relations:} An area of a segment of a sphere with radius $r$
between polar angles $\theta_1$ and $\theta_2$ with
$0 \leq \theta_1 \leq \theta_2 \leq \pi$ is
$\Delta A = 2\pi r^2(\cos\theta_1 - \cos\theta_2)$.

\textbf{Picture of equipment}

\begin{center}\includegraphics[width=0.48\linewidth]{figures/EuPhO_2022_Q4_fig1.jpg}\end{center}

Figure 1: Picture of equipment for experimental problems (protractor and black
paper shield not shown).

\textbf{Illuminance table}

\begin{center}
\begin{tabular}{||c||r|r|r||}
\hline
$T$ [K] & \multicolumn{1}{c|}{Red [lx]} & \multicolumn{1}{c|}{Green [lx]} & \multicolumn{1}{c||}{Blue [lx]} \\
\hline
1570 & 2 & 0 & 0 \\
1600 & 4 & 0 & 0 \\
1610 & 5 & 1 & 0 \\
1620 & 6 & 2 & 0 \\
1630 & 8 & 3 & 0 \\
1640 & 10 & 4 & 0 \\
1660 & 12 & 5 & 0 \\
1670 & 14 & 6 & 0 \\
1700 & 18 & 9 & 1 \\
1730 & 24 & 14 & 3 \\
1780 & 37 & 23 & 7 \\
1820 & 51 & 34 & 11 \\
1880 & 80 & 57 & 21 \\
1940 & 120 & 91 & 36 \\
2000 & 165 & 130 & 53 \\
2060 & 230 & 194 & 80 \\
2120 & 310 & 274 & 118 \\
2160 & 379 & 348 & 155 \\
2220 & 484 & 460 & 210 \\
2260 & 586 & 570 & 264 \\
2310 & 753 & 748 & 348 \\
2350 & 888 & 929 & 440 \\
2390 & 1032 & 1107 & 527 \\
2460 & 1292 & 1452 & 697 \\
2500 & 1577 & 1826 & 879 \\
2540 & 1811 & 2198 & 1078 \\
\hline
\end{tabular}
\end{center}

Table 1: Illuminances by an incandescent light source of a known temperature,
measured through three colour filters at a fixed position of the light source
and the light meter. The accuracy of the measurements is $\pm 2\,\mathrm{lx}$.

\textbf{Power supply instructions}

The benchtop power supply outputs the maximum amount of electricity allowed by
the set current limit and voltage limit. If the current limit is reached first,
it operates as a constant current source, and if the voltage limit is reached
first, it operates as a constant voltage source, indicated by CC and CV
indicator lights.

Do not plug in the light source before setting up the power supply! Notice the
red wire in the image not connected, so the circuit is interrupted.

\begin{center}\includegraphics[width=0.56\linewidth]{figures/EuPhO_2022_Q4_fig2.jpg}\end{center}

You will see two displays: current display on the left and voltage display on
the right. They show the actual current and voltage at the moment, so the
current will be zero, and the voltage will equal the chosen voltage limit. The
CV indicator next to the voltage display will glow, indicating that the voltage
limit is reached.

\begin{center}\includegraphics[width=0.56\linewidth]{figures/EuPhO_2022_Q4_fig3.jpg}\end{center}

Press the voltage knob to set the voltage limit. The display will now show the
voltage limit instead of the actual voltage, and the light below one of the
voltage digits will glow. Turning the knob sets the chosen digit, and pressing
it will cycle through all the digits. Use this to set the voltage to the
maximum allowed voltage of the light source, e.g. $3.00\,\mathrm{V}$ for the
LED.

After 4 seconds, the display goes back to showing the actual voltage and the
light below the voltage digits turns off.

Repeat the procedure with the current knob (left), to set the current limit.

\begin{center}\includegraphics[width=0.56\linewidth]{figures/EuPhO_2022_Q4_fig4.jpg}\end{center}

The example shows the limit $0.200\,\mathrm{A}$. After 4 seconds, the display
reverts to showing the actual current (still zero), and the digit selector
light will turn off.

Now that both limits are set correctly, plug in the light source. Now the
current limit is reached, indicated by the CC (constant current) light. Actual
voltage is now lower than the setting. The light source in this example is not
the same as in your experiment, so your values will be different. Try
increasing the current limit with the left knob.

\begin{center}\includegraphics[width=0.56\linewidth]{figures/EuPhO_2022_Q4_fig5.jpg}\end{center}

When the current is increased too much, the voltage limit is reached (CV
indicator), which protects the light source. Use the current setting to control
the light sources.
